\documentclass{article}


\usepackage[T1]{fontenc}
\usepackage[round]{natbib}
\usepackage{todonotes}
\usepackage{gensymb}
\usepackage{amssymb}
\usepackage{amsmath}
\usepackage{mathtools}
\usepackage[top=1in,bottom=1in,left=1in,right=1in]{geometry}
\usepackage{alltt}
\usepackage{framed}
\usepackage{xcolor}
\usepackage{comment}
\usepackage{graphicx}

\usepackage{xr}[]
\externaldocument{main}
%\usepackage{url}
\usepackage{hyperref}

\renewcommand{\thetable}{\Alph{table}}
\renewcommand{\thefigure}{\Alph{figure}}

\newcommand{\beq}{\begin{equation}}
\newcommand{\eeq}{\end{equation}}
\newcommand{\beqs}{\begin{equation}}
\newcommand{\eeqs}{\end{equation}}
\newcommand{\mx}[1]{\ensuremath{\boldsymbol{\mathrm{{#1}}}}}   % Simple matrix format.


% \newcommand{\review}[1]{\noindent{{\textbf{\textit{#1}}}}\vspace{3mm}}
% \definecolor{colframe}{rgb}{0.01,0.199,0.1}

\newcounter{rtaskno}
\newcommand{\rtask}[1]{\refstepcounter{rtaskno}\label{#1}}

\newcommand{\reviewc}[2]{\begin{leftbar}\noindent\rtask{#1}\ref{#1}. #2\end{leftbar}}
\newcommand{\review}[1]{\begin{leftbar}\noindent #1\end{leftbar}}
\newcommand{\answer}[1]{\noindent #1\\[0mm]}
\newcommand{\copyquote}[1]{\vspace{2mm} {\color{blue} #1}\vspace{2mm}}

\newlength{\leftbarwidth}
\setlength{\leftbarwidth}{3pt}
\newlength{\leftbarsep}
\setlength{\leftbarsep}{10pt}
\renewenvironment{leftbar}{%
    \def\FrameCommand{{\color{black}{\vrule width \leftbarwidth\relax\hspace {\leftbarsep}}}}%
    \MakeFramed {\advance \hsize -\width \FrameRestore }%
}{\endMakeFramed}


% \usepackage[colorlinks=True,citecolor=black]{hyperref}

\begin{document}


\noindent{\bf Poli\v{c}ar and Zupan: 
\\
VERA: Generating Visual Explanations of Two-Dimensional Embeddings \\ via Region Annotation}

\vspace{5mm}
\noindent {\bf COVER LETTER}

\vspace{5mm}


\noindent Dear Editor and the Reviewers, \\
 
\vspace{3mm}

\noindent We would like to submit our manuscript ``\textit{VERA: Generating Visual Explanations of Two-Dimensional Embeddings via Region Annotation}'' to be considered for publication in the ``\textit{IEEE Transactions on Visualization and Computer Graphics}''.
We had submitted an earlier version of this manuscript to last year's IEEE Vis conference, where the reviewers raised several valid concerns that we have addressed in our revised paper. We would like to thank these reviewers for their valuable, detailed feedback, which helped us to improve the manuscript.

Below we provide point-by-point responses to the comments and concerns raised by the reviewers.

\vspace{3mm}

\noindent Kind regards,
\vspace{3mm}

\noindent Pavlin G. Poli\v{c}ar,
Bla\v{z} Zupan \\ 
\noindent Ljubljana, XX, 2025

\vspace{10mm}


\subsection*{Reviewer 1}

\reviewc{r1q1}{W1: From my understanding of the authors' system, I disagree with their base assertion that the system does not require some degree of human intervention.  To begin the process, the embedding algorithm itself needs to be selected by a user, who must also tune the hyperparameters to produce an embedding that the VERA system can ingest.  It can be fairly argued that this step is external to the VERA automated process, but it is undeniable that the system is dependent on receiving an embedding that contains meaningful regions, a fact that is noted by the authors in the final paragraph of Section 3.1.2.}

\answer{}


\reviewc{r1q2}{W2: Even if the previous point is ignored, the authors themselves state a number of human-tunable parameters that are necessary to adjust to that input embedding to produce a high-quality output from the VERA system, including level of KDE-based contour lines (Section 3.1.2), mapping distortion estimates for each data point (also 3.1.2), and region overlap threshold to determine whether regions should be joined (Section 3.1.3).  It could be argued that each of these parameters has a default and that setting them is optional, but using the defaults does not appear to guarantee a quality, meaningful VERA output.}

\answer{}


\reviewc{r1q3}{W3: While the embedding in Figure 2 that is used throughout the description of the system (Section 3), I found it to be very artificial in both the embedding layout, and the simplicity of the relationships between input features (e.g., note that Country Purchased and Number of Children have identical regions).  Using a more complex data set, such as that provided in the teaser figure and Section 4 would better emphasize the power of the VERA technique (or would further highlight some of the limitations).}

\answer{While we agree that the guiding example used throughout the paper is very simple, we designed this dataset specifically in this way so that the different aspects and steps of our procedure could be demonstrated in a clear and visual manner. Moreover, while simple, this data set reflects many of the challenges relevant to our method in numerous real-world data sets.  However, we agree that including more real-world examples would better demonstrate the capabilities of VERA, therefore, we have added several more complex, real-world data sets to Section~\ref{sec:case-studies}.}


\reviewc{r1q4}{W4: On the limitations point, the authors do note that the system is limited to datasets in which features have "human-readable, semantic meaning" and also encounters scaling problems as the number of features increases.  I would add several more limitations to these.  (1) the process of discretizing numeric variables into five bins assumes a normal of uniform distribution, whereas other distributions (skewed, bimodal) might be separated in a better way; (2) unless the term "binary indicator variable" is being incorrect used, these five bins still require further breakdown into five separate variables (as do any categorical variables with >2 classes), making the technique effectively limited to binary classification rather than naturally supporting multi-class classification;  (3) the process of merging variables assumes that consecutive order is important (x<5, 5<x<6, x>6), but there are cases such as correlations where values close to -1 and 1 share more in common than those close to 0.}

\answer{}


\reviewc{r1q5}{M2:  The system makes use of both descriptive and contrastive explanations, but only contrastive explanations are clearly called out in the Related Work.  Adding the "description explanation" keywords to the appropriate paragraphs in Section 2 will better highlight the differences between VERA and existing systems.}

\answer{}


\reviewc{r1q6}{M3:  At the end of Section 3.1.3, it is noted that "Features whose values exhibit no spatial correlation with the embedding space are merged into a single explanatory variable."  What percentage of features fall under this category in a typical embedding?  And does this number decrease if users tune the parameters noted in W1 to create a better embedding or in W2 to alter how the embedding is parsed?}

\answer{}


\reviewc{r1q7}{M4:  In Figure 5, the Language plot displays English in two distinct regions of the embedding rather than merging them into a single joined region.  Does this separation have an impact on user understanding, particularly after multiple features are joined to define a region?}

\answer{}


\reviewc{r1q8}{M5:  Was there any thought given to extending explanations to describe why regions are (or are not) joined?  It's outside the scope of this system, but I'm thinking of the MNIST Digits dataset in particular.  How would a user interpret the decision underlying joining similar numbers like 1 and 7 but not other similar numbers like 6 and 8?}

\answer{}


\reviewc{r1q9}{M6:  In the results in Section 5.2, reporting standard deviations along with the means and p-values would be useful context.}

\answer{Agreed. We have now included measures of standard deviations and standard errors where appropriate.}


\subsection*{Reviewer 2}

\reviewc{r2q1}{While the paper overall is well written and motivated, I cannot quite connect the
dots for one of the main story points. Understanding embedding projections is
challenging, yes. So this paper proposes a way to solve that, but through a static
visualization. But the contribution is primarily an algorithm to label regions of
a projection.}

\answer{}


\reviewc{r2q2}{The evaluation that was added to the end of the paper after some
discussion from the previous reviews then compares the static visualization to an
interactive visualization of the same thing, but without the automatic VERA
labeling. This evaluation seems odd to me, as my expected comparison would have
been how this algorithm compares to other projection labeling algorithms. I
totally get the synergy of combining VERA with an interactive tool - it seems like
every interactive tool should have the VERA approach built-in and be the default
"on load" screen!}

\answer{}


\reviewc{r2q3}{Given the length of the algorithm description, I expected the
evaluation to be on it compared to other labeling algorithms, either in runtime
performance, quality, or measuring user understanding. This relates to my next
point below.}

\answer{}


\reviewc{r2q4}{Previous reviews ask about how the approaches scales. There has been some new
texted added explaining that as features grow to the hundreds of thousands this
approach may not work well since the labels will be long. It proposes using LLMs
to summarize, which is neat, but my concerns are moreso for how this approach
scales not just in features but also in number data points. For example, common
use of dimensionality reduction is for ML datasets, which can have millions of
points. There is no mention of scale in the current submission, and the only
examples we see appear to have just a few hundred points.

Moreover, I have other questions about scalability that isn't touched upon in the
paper, which I found odd given it's strong algorithmic contribution. For example,
how fast does this take to run? Is it in minutes, seconds, or interactive speed?
How was it implemented? Where can it run? I see the code link but understand it's
private until acceptance.}

\answer{}


\reviewc{r2q5}{While using dimensionality reduction of larger, labelled, high-dimensional dataset
is useful for certain domains (mentioned in the submission), another common
application for using this to cluster data for exploration in a unsupervised
learning approach. Oftentimes in this case there is no additional
metadata/features, or at least none that have interpretable units. It is written
off that VERA would need to be extended to tackle this. Is there anything more
that could be said here? Any hints, ideas, conjectures of applying VERA to these
datasets?}

\answer{}


\reviewc{r2q6}{I greatly appreciate the detailed and natural language algorithm description, as
it walks a reader through the different types of explanations. Using a synthetic
dataset with known features and clusters is a good way to compare results (the
figures emphasize this). I will say the amount of new terminology introduced is a
little overwhelming at first since it's multiple levels of groups of points (lots
of aggregating and merging!). Despite these pros, this section is rather long.
Nearly 50\% of the paper is dedicated to this section and description. This balance
makes the remainder of the paper feel rushed, and doesn't leave room for more
examples, which I explain below.}

\answer{}


\reviewc{r2q7}{Since a major contribution of this work is the algorithm and the static labelled
projection visualizations, it would be great to see more results on different
datasets. This was also highlighted in the previous reviews, and while the new
user study includes more than just the synthetic books datasets, we still don't
get to see many many visualizations. I'd argue that showing many more examples
would strongly help this submission given it's emphasis on static representation.
I imagine a new figure (could even be full page, however I understand this would
require cutting content - referring to my other concern above the algorithm
sections could be compressed) that shows 3, 5, or even a dozen different
embeddings as small multiples with the VERA approach applied to them to give
readers a better idea of the approach and results.}

\answer{}


\reviewc{r2q8}{It would be nice to have the related work partitioned, say into subsections or
using paragraph titles to provide short descriptions for the different topics.
While the related work is good, having these would help readability.}

\answer{}



\subsection*{Reviewer 3}

\reviewc{r3q1}{W1: The introduction could be strengthened by highlighting the techniques in VERA
(e.g., contrastive and descriptive explanation selection/ranking) and user study
results. While the current introduction provides background and motivation for
embedding explanations, it lacks clarity on how VERA addresses mentioned
challenges and its effectiveness.}

\answer{}


\reviewc{r3q2}{W2: The related work is comprehensive but overlooks recent interactive embedding
visualization tools [e.g., 1–3]. For instance, WizMap [1] can automatically
generate interactive embedding explanations. How does VERA differ from this work?
\vspace{3mm}
\\
\noindent [1] Wang, Z. J., Hohman, F., \& Chau, D. H. (2023). WizMap: Scalable Interactive
Visualization for Exploring Large Machine Learning Embeddings. arXiv preprint
arXiv:2306.09328.
\\
\noindent [2] Lekschas, F. (2023). Regl-scatterplot: A scalable interactive javascript-based
scatter plot library. Journal of Open Source Software, 8(84), 5275.
\\
\noindent [3] Sivaraman, V., Wu, Y., \& Perer, A. (2022, March). Emblaze: Illuminating
machine learning representations through interactive comparison of embedding
spaces. In 27th International Conference on Intelligent User Interfaces (pp.
418-432).}

\answer{Thank you for suggesting the additional references. We have now included a citation to WizMap, which is indeed relevant to our work. We now also mention three approaches that propose methods to compare multiple embeddings in the introductory paragraph of the Related Work section, including a citation to Emblaze. While not directly related to our work, comparing embeddings to one another is an adjacent task to ours. We have decided not to include the citation to Regl-scatterplot, since the publication describes a specific Javascript scatter plot rendering library which is not used in our work.}


\reviewc{r3q3}{W3: VERA assumes that the features in the high dimensional space are
interpretable. Many people use tSNE and UMAP to generate 2D embedding
visualizations for neural network's latent representation embeddings. These
embeddings have a large number of dimensions (e.g., 1024 dimensions), and each
dimension does not map to a human-understandable feature. How does VERA generalize
to these ML embeddings? This is briefly discussed in the limitation section. I
recommend mentioning the focus on tabular/interpretable features earlier in the
paper.}

\answer{}


\reviewc{r3q4}{W4: Section 3 mentions some configurations in VERA, such as k-means quantization
with k=5 and using top-3 explanation selection. How can users adjust these
settings in practice? This is not detailed in Sections 4 and 5. It would be great
if VERA is more interactive and allows users to easily change configurations
through interactions (e.g., slider, drop-down menu).}

\answer{}


\reviewc{r3q5}{W5: Using Orange data mining toolkit as a baseline model may not be ideal since it
lacks embedding explanations. I wonder if it is possible to use related works with
embedding explanations as a stronger baseline model. Otherwise, it should be
acknowledged in the limitation section.}

\answer{}


\reviewc{r3q6}{M1: What does "satisfying explanations" mean? (Section 3) One might argue that
"satisfying explanations" can elicit confirmation bias when interacting with VERA.}

\answer{}


%\bibliographystyle{abbrvnat}
%\bibliography{main}

\end{document}
